\documentclass[10pt,a4paper]{amsart}
\usepackage[margin=19mm]{geometry}
\usepackage{amsmath,amssymb,mathtools}
\usepackage[scr=rsfs,cal=euler]{mathalfa}
\usepackage{xfrac}
\usepackage[unicode,hidelinks,bookmarks=false]{hyperref}
\newcommand{\ie}{i.e.\ }
\newcommand{\cf}{cf.\ }
\newcommand{\resp}{respectively\ }
\newcommand{\Subsection}{\subsection}
\providecommand{\nobreakdash}{\nobreak\hbox{-}}
\setlength{\emergencystretch}{2em}
\multlinegap0pt
\usepackage{bm}
\usepackage{bbm}
\usepackage{dsfont}
\usepackage{xspace}
\usepackage{cancel}
\usepackage{mathtools}
\usepackage{amsthm}
\makeatletter
\@ifundefined{theorem}{\newtheorem{theorem}{Theorem}}{}
\@ifundefined{lemma}{\newtheorem{lemma}[theorem]{Lemma}}{}
\makeatother
\title[Hodge structures of the surface of planes in a cubic 5-fold]{Hodge structures of the surface of planes in a cubic 5-fold}
\author{Chenpeng Feng}
\date{Source: arXiv:2506.13767v1 (CC BY 4.0)}
\begin{document}
\maketitle

\noindent\emph{Source and licence.} Hodge structures of the surface of planes in a cubic 5-fold. Version arXiv:2506.13767v1 (CC BY 4.0), distributed under the Creative Commons Attribution 4.0 International licence, as recorded in the arXiv licence metadata for this exact version and shown on the abstract page https://arxiv.org/abs/2506.13767v1. Full rights notice: licenses/arxiv-2506.13767-CC-BY-4.0.txt.\par\smallskip

\noindent\textit{Editorial scope.} Counted unit: the Lemma whose source label is \texttt{LmmChernClassesOfSurfaceInP3} in the article (source0.tex lines 564--570, source label \texttt{LmmChernClassesOfSurfaceInP3}), delivered together with the article's own complete original proof of it (source0.tex lines 571--604, one \texttt{proof} environment of 34 lines). No mathematics is written, repaired or re-derived: statement and proof are the article's own text line for line. Required context retained uncounted: none. Every definition, display and label of those lines is retained unchanged. Omitted from this package and not redistributed: 1--563, 605--1247. Nothing inside a retained excerpt is omitted or altered: every retained body is delivered exactly as the article's lines are written, so no omission receipt of any kind accompanies this package. Citations: the retained excerpts carry no citation command at all, so no bibliography entry of the article is transcribed and no reference is redistributed. Reference adaptation: none was needed and none was made; every reference inside the delivered unit resolves inside it, so no unresolved reference or citation is produced. Printed slips kept literal: no printed word, symbol, hypothesis or number of the counted statement or of its proof was altered. Layout and class: the article's own class is \texttt{article} with packages including amsmath, amssymb, amsthm, mathtools, nicematrix, hyperref, titlesec, tikz-cd, geometry, graphics, setspace; the wrapper is the lane's pinned \texttt{amsart} editorial wrapper at 10pt on A4, which restates the theorem-like environments the retained text needs and adds the article's own macro definitions that the retained text needs (none), with extra wrapper packages (bm, bbm, dsfont, xspace, cancel, mathtools). The delivered numbering is the unit's own. Assets: no retained span contains a figure environment, a graphics-inclusion command, a PDF-page inclusion, a table or a TikZ picture; no image file is copied into this package, the package declares no asset dependency and no image file is redistributed with it. Licence: version arXiv:2506.13767v1 of this article is distributed under the Creative Commons Attribution 4.0 International licence, as recorded in the arXiv licence metadata for this exact version and shown on the abstract page https://arxiv.org/abs/2506.13767v1. A full rights notice is delivered at licenses/arxiv-2506.13767-CC-BY-4.0.txt. Characters: every retained excerpt is pure ASCII, so the delivered engine, XeLaTeX with its default Latin Modern text font, reports no missing character.
\begin{lemma}\label{LmmChernClassesOfSurfaceInP3}
For a smooth surface \( S \) of degree \( d \) in \(\mathbb{P}^3\),
the first Chern class and the second Chern class of the tangent bundle $T_S$ over $S$ are:
\[ c_1(S) = (4-d)H. \]
\[ c_2(S) = d^3 - 4d^2 + 6d. \]
Moreover, we have $c_1(S)^2 = d^3 - 8d^2 + 16d$.
\end{lemma}

\begin{proof} 
Consider the short exact sequence
\[ 
0 \to T_S \to T_{\mathbb{P}^3}|_S \to N_{S|\mathbb{P}^3} \to 0. 
\]
Here, \(T_S\) is the tangent bundle of \( S \), \( T_{\mathbb{P}^3}|_S \) is the restriction of the tangent bundle $T_{\mathbb{P}^3}$ of \(\mathbb{P}^3\) to \( S \), and \( N_{S|\mathbb{P}^3} \) is the normal bundle of \( S \) in \(\mathbb{P}^3\). From the above sequence, we have 
$$c(T_{\mathbb{P}^3}|_S) = c(S) \cdot c(N_{S|\mathbb{P}^3}).$$ 
We know that the total Chern class of \( T_{\mathbb{P}^3} \) is $c(T_{\mathbb{P}^3}) = (1 + h)^4$,
where \( h \) is the hyperplane class in \(\mathbb{P}^3\). Restricting to \( S \), we get 
$$c(T_{\mathbb{P}^3}|_S) = (1 + h|_S)^4=(1+H)^4$$ 
where $H:= h|_S $.  

Now, the normal bundle \( N_{S/\mathbb{P}^3} \) is \(\mathcal{O}(d) \), so its total Chern class is
\[ c(N_{S|\mathbb{P}^3}) = 1 + dH. \]
Hence we have 
\[ (1 + H)^4 = c(S) \cdot (1 + dH). \]
Expanding and comparing terms, we get
\[
c_1(S) = (4 - d)H, 
\]
\[
c_2(S)  = (6 - 4d + d^2)H^2. 
\]

Hence  \[ 
c_1(S)^2 = ((4 - d)H)^2 = (4 - d)^2H^2 = (16 - 8d + d^2)H^2.
\]
Since \( H^2 \) represents the self-intersection of two hyperplane sections on a smooth surface of degree \( d \) in \(\mathbb{P}^3\), we have \( H^2 = d \). Therefore we have
$$ 
c_1(S)^2 =d^3 - 8d^2 + 16d.
$$
and
$$c_2(S) = d^3 - 4d^2 + 6d. $$
\end{proof}

\end{document}
